% --- преамбула для "матфиз" стиля (pdfLaTeX) ---

% Файл: Лекция6_верстка.tex
% Сверстано по твоему референсу (малые заголовки, строгая последовательность окружений)
\documentclass[12pt,a4paper]{article}

\usepackage[utf8]{inputenc}
\usepackage[T2A]{fontenc}
\usepackage[russian]{babel}

\usepackage{amsmath,amssymb,amsthm}
\usepackage{mathtools}
\usepackage{geometry}
\usepackage{microtype}
\usepackage{enumitem}
\usepackage{graphicx}
\usepackage{caption}
\usepackage{hyperref}
\usepackage{xcolor}

\geometry{left=25mm,right=25mm,top=25mm,bottom=25mm}

\usepackage{titlesec}
\titleformat{\subsubsection}{\normalfont\normalsize\bfseries}{\thesubsubsection}{1em}{}
\titlespacing*{\subsubsection}{0pt}{10pt}{4pt}
\titleformat{\paragraph}[runin]{\normalfont\normalsize\bfseries}{}{}{}[\quad]
\newcommand{\qheading}[1]{\paragraph{#1}}

\theoremstyle{definition}
\newtheorem{definition}{Определение}[subsubsection]

\theoremstyle{plain}
\newtheorem{theorem}[definition]{Теорема}
\newtheorem{lemma}[definition]{Лемма}
\newtheorem{corollary}[definition]{Следствие}

\theoremstyle{remark}
\newtheorem{remark}[definition]{Замечание}
\newtheorem{example}[definition]{Пример}

\renewenvironment{proof}[1][Доказательство]{%
	\par\noindent\textit{#1.}\quad
}{\hfill\(\blacksquare\)\par}

\numberwithin{equation}{subsubsection}
\captionsetup{font=small,labelfont=bf}
\setlist[itemize]{noitemsep,topsep=4pt,parsep=0pt,partopsep=0pt}

\begin{document}
	
	\begin{center}
		{\Large\bf Лекция 6. Постановка, существование и единственность начально-краевых задач для уравнений гиперболического типа}
	\end{center}
	
	\subsubsection{Постановка начально-краевой задачи}
	Пусть \(B\subset\mathbb{R}^n\) — область с гладкой границей \(\partial B\). Рассматриваем уравнение гиперболического типа (в общем виде)
	\[
	\text{(PDE)}\qquad U_{tt} = \mathcal{L} U + F(t,x),
	\]
	где \(\mathcal{L}\) — второй пространственный дифференциальный оператор (например, оператор Лапласа с коэффициентами), \(x\in B\), \(t>0\), и задано граничное условие на \(\partial B\) и начальные условия при \(t=0\):
	\[
	U|_{t=0}=U_0(x),\qquad U_t|_{t=0}=V_0(x).
	\]
	(В оригинале лекции формулировка даётся в конкретном виде — см. источник.) :contentReference[oaicite:2]{index=2}
	
	Граничные условия могут быть различных типов:
	\begin{enumerate}
		\item Дирихле: \(U|_{\partial B}=f(t,\xi)\).
		\item Нейман: \(\partial_n U|_{\partial B}=g(t,\xi)\).
		\item Смешанные (например, Робина и т.п.).
	\end{enumerate}
	
	\subsubsection{Условие согласования}
	Необходимы условия согласования начальных и граничных данных (чтобы решение было непрерывно до границы и в момент \(t=0\)). В лекции это перечислено в виде трёх требований на допустимые множества и значения данных. :contentReference[oaicite:3]{index=3}
	
	\subsubsection{Единственность решения (схема доказательства)}
	\begin{theorem}
		(Единственность решения начально-краевой задачи)  
		Пусть две функции \(U^{(1)}\) и \(U^{(2)}\) удовлетворяют одной и той же задаче (одинаковые уравнение, граничные и начальные условия) и обладают необходимой регулярностью. Тогда \(U^{(1)}\equiv U^{(2)}\) в области рассмотрения.
	\end{theorem}
	
	\begin{proof}
		Пусть \(W=U^{(1)}-U^{(2)}\). Тогда \(W\) удовлетворяет однородной задаче (с нулевыми начальными и граничными данными):
		\[
		W_{tt} = \mathcal{L} W,\qquad W|_{t=0}=0,\ W_t|_{t=0}=0,\ W|_{\partial B}=0.
		\]
		Введём подходящую энергетическую функционалу (энергию)
		\[
		E(t)=\frac12\int_B \big( |W_t|^2 + a(x)\,|\nabla W|^2 \big)\,dx,
		\]
		и получим интегральное тождество (энергетическое) вида
		\[
		\frac{d}{dt}E(t) = \text{(граничные и нулевые члены)}.
		\]
		Интегрирование по времени от \(0\) до \(t\) даёт \(E(t)=0\) при всех \(t\), откуда \(W\equiv 0\). Это завершает доказательство единственности. :contentReference[oaicite:4]{index=4}
	\end{proof}
	
	\begin{corollary}
		Следует, что при выполнении условий согласования решение, если существует, единственно.
	\end{corollary}
	
	\qheading{Замечание}
	В исходной лекции даётся более подробное интегральное тождество с отображением площадных и объёмных интегралов; я вынес сюда только ключевую идею (энергетический метод). См. источник для подробной формы тождеств. :contentReference[oaicite:5]{index=5}
	
	\subsubsection{Интегральное тождество (пример)}
	В лекции приводится интегральное тождество, получаемое при умножении уравнения на тестовую функцию \(V\) и интегрировании по области \(B\):
	\[
	\int_B (U_{tt}V + \nabla U\cdot\nabla V)\,dx = \int_B F V\,dx + \int_{\partial B}(\cdots)\,dS,
	\]
	и дальнейшие преобразования позволяют вывести оценку энергии. (Полная формула в исходнике.) :contentReference[oaicite:6]{index=6}
	
	\subsubsection{Метод Даламбера (задача о колебаниях струны)}
	Далее в лекции рассматривается применение общего интеграла (метода Даламбера) для уравнения волны на прямой:
	\[
	u_{tt}-c^2 u_{xx}=0,\qquad x\in\mathbb{R},\ t>0,
	\]
	с начальными условиями
	\[
	u(x,0)=\phi(x),\qquad u_t(x,0)=\psi(x).
	\]
	Общее решение (Даламбер) имеет вид
	\[
	u(x,t)=\frac{\phi(x-ct)+\phi(x+ct)}{2}+\frac{1}{2c}\int_{x-ct}^{x+ct}\psi(s)\,ds.
	\]
	
	\qheading{Частные случаи и интерпретация}
	Если начальное возмущение финитно (имеет компактный носитель), то с течением времени возмущение расползётся вдоль характеристик; описывается появление переднего и заднего фронтов волн (см. рисунки в исходнике). Также обсуждается случай конечной струны длины \(l\) с закреплёнными концами и явление отражений.
	
	\subsubsection{Замечание об устойчивости решения}
	В лекции подчеркивается устойчивость решения по правым частям: если правая часть уравнения \(F\) мала в некоторой норме, то и соответствующая компонента решения имеет такой же порядок малости (оценки в энергетических нормах).
	
	% ----- конец лекции 6 -----
	\vspace{1em}
	
	$\dfrac{1}{\pi} \displaystyle\int\limits_{0}^{+\infty} \dfrac{x - \sin x}{x^3} dx$
	
\end{document}


